%% kapa: exam=1 round=28 year=2017 period=1 form=A subject=real_estate subject_ko=부동산학원론 seq=17 number=97
%% source: Q-Net 감정평가사 기출문제 (cst00309) artlSeq=5199603 file=extracted/1008/5199603/2017년 제28회 감정평가사 1차시험 문제지/1교시 A형.pdf
%% conversion: kordoc 4.24.51 -> tools/gen1.py (auto); review=proofread-sonnet; flags=none
\begin{question}{97}
토지의 정착물과 동산에 관한 설명으로 옳지 않은 것은?
\begin{choices}
  \item 부동산과 동산은 공시방법을 달리하며, 동산은 공신의 원칙이 인정되나 부동산은 공신의 원칙이 인정되지 않는다.
  \item 토지의 정착물 중 명인방법을 구비한 수목의 집단은 토지와 독립적인 거래의 객체가 될 수 있다.
  \item 토지의 정착물 중 도로와 교량 등은 토지와 독립적인 것이 아니라 토지의 일부로 간주된다.
  \item 제거하여도 건물의 기능 및 효용의 손실이 없는 부착된 물건은 일반적으로 동산으로 취급한다.
  \item 임차인이 설치한 영업용 선반ㆍ카운터 등 사업이나 생활의 편의를 위해 설치한 정착물은 일반적으로 부동산으로 취급한다.
\end{choices}
\end{question}
\answer{5}
